\documentclass[11pt,a4paper]{exam}

\usepackage[utf8]{inputenc}
\usepackage[T1]{fontenc}
\usepackage[english,main=french]{babel}
\usepackage{amsmath, amssymb}
\usepackage{geometry}
\usepackage{listings}
\usepackage{xcolor}

\geometry{margin=2cm}

\lstset{
  basicstyle=\ttfamily\small,
  keywordstyle=\bfseries,
  columns=fullflexible,
  keepspaces=true,
  frame=single,
  showstringspaces=false,
  tabsize=2
}

\firstpageheader{BTS CIEL1}{Python}{Date : 02/12/2025}
\runningheader{Python}{}{Page \thepage/\numpages}

\begin{document}

\begin{center}
  \Large\textbf{Algorithmes et structures de données en Python}\\[0.5em]
  \normalsize Durée : 1h \\
  \vspace{1em}
  \textbf{Nom :} \rule{6cm}{0.4pt} \hfill
  \textbf{Prénom :} \rule{6cm}{0.4pt}\\[1em]
\end{center}

\vspace{1em}

\section*{Exercice 1 — QCM (10 questions)}

Cochez la bonne réponse pour chaque question. \textbf{Une seule} réponse correcte par question.

\begin{questions}

\question[1] Quel type de langage est Python ?
\begin{checkboxes}
  \choice Un langage compilé
  \choice Un langage interprété
  \choice Un langage matériel
  \choice Un langage dépendant du système
\end{checkboxes}

\question[1] Dans une liste Python, que permet \texttt{append()} ?
\begin{checkboxes}
  \choice Trier la liste
  \choice Supprimer un élément
  \choice Ajouter un élément à la fin
  \choice Copier la liste
\end{checkboxes}

\question[1] Une liste se différencie d'un tuple car :
\begin{checkboxes}
  \choice Elle est immuable
  \choice Elle peut être modifiée
  \choice Elle contient un seul type d’élément
  \choice Elle ne peut pas être parcourue en boucle
\end{checkboxes}

\question[1] Dans un dictionnaire, la condition \texttt{"nom" in d} vérifie :
\begin{checkboxes}
  \choice Si la valeur existe
  \choice Si la clé existe
  \choice Si le dictionnaire est vide
  \choice Si la longueur du dictionnaire est supérieure à 1
\end{checkboxes}

\question[1] Qu’est-ce qu’un élément d’un ensemble (set) ?
\begin{checkboxes}
  \choice Un élément identifié par un indice
  \choice Un élément qui peut être dupliqué
  \choice Un élément unique, sans ordre particulier
  \choice Une paire clé-valeur
\end{checkboxes}

\question[1] Quelle valeur est considérée comme \texttt{False} en Python ?
\begin{checkboxes}
  \choice 10
  \choice \texttt{''} (chaîne vide)
  \choice \texttt{"Hello"}
  \choice 3.14
\end{checkboxes}

\question[1] La boucle \texttt{for} en Python permet :
\begin{checkboxes}
  \choice D’exécuter un code tant qu’une condition est vraie
  \choice De parcourir les éléments d’un objet itérable
  \choice De définir une fonction
  \choice De gérer les erreurs
\end{checkboxes}
\newpage

\question[1] La fonction \texttt{range(2, 7)} génère :
\begin{checkboxes}
  \choice 2, 3, 4, 5, 6
  \choice 2, 3, 4, 5, 6, 7
  \choice 3, 4, 5, 6, 7
  \choice 2, 7
\end{checkboxes}

\question[1] Parmi les opérations ci-dessous, laquelle modifie une liste en place ?
\begin{checkboxes}
  \choice \texttt{sorted(lst)}
  \choice \texttt{lst.copy()}
  \choice \texttt{lst.sort()}
  \choice \texttt{lst + lst}
\end{checkboxes}

\question[1] Laquelle de ces structures est la plus adaptée pour associer un prix à chaque produit ?
\begin{checkboxes}
  \choice Liste
  \choice Tuple
  \choice Ensemble
  \choice Dictionnaire
\end{checkboxes}

\end{questions}

\section*{Exercice 2 — Programme Python à trous}

\begin {questions}

\question[2] Complétez le programme Python (Micro:bit) ci-dessous :

\begin{lstlisting}[language=Python]
from microbit import *

import ________

MAX = ________

number_to_guess = random.________(1, MAX)
answer = 0

display.________()

while True:
    display.show(question_mark_image)
            
    if ________.is_pressed():
        answer = (answer + 1) % (MAX + 1)
        display.scroll(answer, delay=50, wait=True)

    if ________.was_pressed():
        if answer > number_to_guess:
            display.show(Image.SAD)
            display.show(Image.LESS)
            answer = 0
        if answer < number_to_guess:
            display.show(Image.SAD)
            display.show(Image.PLUS)
        else:
            display.show(Image.HAPPY)
            answer = 0
        sleep(_______)

\end{lstlisting}

\question[0.5] Expliquez l'instruction suivante : \verb|(answer + 1) % (MAX + 1)|

\end {questions}

\newpage

\section*{Exercice 3 — Traduction de C vers Python}

\begin{questions}
\question On donne le programme C suivant :

\begin{lstlisting}[language=C]
#include <stdio.h>

int main(int argc, char const *argv[]) {
  int e = 0;
  int nbOcc = 0;
  int posOcc = -1;
  int n = -1;
  int tab[100] = {0};

  printf("Combien d'elements souhaitez-vous saisir :");
  scanf("%d", &n);

  for (int i = 0; i < n; i++) scanf("%d", &tab[i]);

  printf("Quel element souhaitez-vous rechercher :");
  scanf("%d", &e);

  for (int i = 0; i < n; i++) {
      if (tab[i] == e) {
          nbOcc++;
          posOcc = i;
      }
  }
  printf("Nombre d'occurences : %d, position : %d \n", nbOcc, posOcc);
}
\end{lstlisting}

\begin{parts}
  \part[0.5] Expliquez en quelques mots ce que fait ce programme.
  \vspace{3cm}

  \part[2] Transcrivez ce programme en langage \textbf{Python} ci-dessous :

  \vspace{0.5cm}
  \begin{solutionbox}{8cm}
  \end{solutionbox}
\end{parts}
\end{questions}

\newpage

\section*{Exercice 4 — Programme Micro:bit}

La carte \textbf{micro:bit} possède une \textbf{matrice de 25 LED} (5 colonnes × 5 lignes).  

Chaque LED est identifiée par ses coordonnées :

\begin{itemize}
    \item \texttt{x} : numéro de la colonne (0 à 4)
    \item \texttt{y} : numéro de la ligne (0 à 4)
\end{itemize}

Pour lire ou modifier l'état d'une LED, on peut utiliser :

\begin{itemize}
    \item \verb|display.get_pixel(x, y)| : renvoie l'état actuel (0 = éteinte, 9 = allumée)
    \item \verb|display.set_pixel(x, y, valeur)| : modifie l'état
\end{itemize}

\begin{questions}

\question[5] \textbf{Ecrivez un programme Python qui :}
\begin{itemize}
    \item génère une liste de nombres aléatoires (entre 0 et 4 inclus),
    \item dessine sur la matrice de LED une colonne:
    \begin{itemize}
        \item l'indice dans liste correspond à position de la colonne (la coordonée x),
        \item la valeur dans la liste correspond à la hauteur de la colonne (la coordonée y)
    \end{itemize}
    \item lors de l'appuie sur le bouton A :
    \begin{itemize}
      \item efface l'écran et re-génére la liste.
    \end{itemize}
\end{itemize}

\vspace{0.5cm}
\noindent\textbf{Votre programme Python :}

\begin{solutionbox}{15cm}
\end{solutionbox}

\end{questions}


\end{document}
